% anexa-latex.tex starts here
\chapter{Anexă: Redactarea lucrării în \LaTeX}\label{ch:appendix-latex}

\section{Instrumente pentru redactare}
Lucrarea a fost redactată cu motorul \texttt{LuaLaTeX}, folosind fonturi Unicode și pachetul
\texttt{biblatex} pentru gestionarea bibliografiei. Structura este modulară, fiecare capitol
fiind un fișier separat (\texttt{.tex}) compilat prin scripturile automate de build.
Imaginile sunt salvate în \texttt{images/}, iar citările se realizează cu \verb|\cite{}|.

\section{Compatibilitate multiplatformă}
Există scripturi automate pentru Linux (\texttt{build.sh}), Windows (\texttt{build.bat})
și PowerShell (\texttt{build.ps1}) care generează PDF-ul în directorul \texttt{dist/}
și afișează eventualele erori de compilare.

\vspace{0.5em}
\noindent\textit{Notă:} Mediul \texttt{TreeVerb} utilizează fontul \texttt{FreeMono} pentru suport complet al caracterelor Unicode de tip box-drawing (U+2500–U+257F).

% ----------------------------------------------------------
% Include non-cited but relevant sources to keep bibliography complete
% ----------------------------------------------------------
% These sources are foundational to the broader PAN architecture
% but were not cited directly in the main chapters (e.g., roadmap tools).
\nocite{vite,cockroachdb,redis,firebase_auth,json_schema}

% --- PAN Architecture Manifest ---

\chapter{Kookio: Manifest Arhitectural PAN (Prisma – Analog – Nx)}\label{ch:appendix-manifest}

\section{Viziune generală}
\textbf{Kookio} este o aplicație full-stack dezvoltată în \texttt{TypeScript}, bazată pe arhitectura
\textbf{PAN} — un cadru scalabil care integrează \textbf{Prisma} pentru acces la date,
\textbf{Analog} pentru randare hibridă și \textbf{Nx} pentru orchestrare monorepo.
Această abordare reprezintă o alternativă la stack-uri consacrate (JAMstack, MEAN),
fiind adaptată ecosistemului \texttt{Angular} și cerințelor enterprise.

\section{Componentele arhitecturii}
\begin{itemize}
  \item \textbf{Prisma} – ORM cu suport pentru PostgreSQL și extensii NewSQL (CockroachDB).
  \item \textbf{Analog} – meta-framework pentru Angular cu SSR, Signals și tRPC.
  \item \textbf{Nx} – orchestrator modular care optimizează build-urile și testele.
\end{itemize}

\section{Structura fizică a proiectului}
Structura completă a monorepo-ului este ilustrată în \textbf{§~\ref{subsec:monorepo-struct}} din capitolul~\ref{ch:impl}.

\section{Avantaje comparative}
\begin{table}[H]
  \centering
  \small
  \begin{tabularx}{\linewidth}{@{} Y Y Y Y @{}}
    \toprule
    \thead{Caracteristică} & \thead{JAMstack} & \thead{MEAN} & \thead{PAN (Kookio)} \\
    \midrule
    SSR / Hybrid Rendering  & Parțial (Next.js) & Nu & Da (Analog SSR) \\
    Type-safety end-to-end  & Parțial           & Nu (Mongo/JS) & Da (TS + Prisma + Angular) \\
    Structurare modulară    & Slabă             & Medie         & Excelentă (Nx + libs) \\
    Acces la date           & Headless CMS/REST & REST/Mongoose & ORM modern (Prisma) \\
    Adaptabilitate monorepo & Nu                & Nu            & Da \\
    \bottomrule
  \end{tabularx}
  \caption{Comparație între arhitectura PAN și alte stack-uri moderne}
  \label{tab:pan-adv}
\end{table}

\section{Direcții de viitor}
Direcțiile viitoare de dezvoltare sunt prezentate în \textbf{§~\ref{sec:roadmap}}, unde sunt discutate extensiile funcționale și tehnice ale aplicației.

\section{Concluzie}
Arhitectura \textbf{PAN} sintetizează cele mai bune practici moderne din ecosistemul \texttt{TypeScript}, oferind o experiență unificată full-stack și o bază solidă pentru aplicații educaționale, alimentare sau de lifestyle, precum viziunea \textit{Kookio}.
\textbf{PAN} nu reprezintă un standard formal, ci un model arhitectural derivat din integrarea unor tehnologii mature, validat prin implementarea aplicației \textit{Kookio}.
% anexa-latex.tex ends here
